\documentclass[11pt]{article}
\usepackage[margin=0.95in]{geometry}
\usepackage[T1]{fontenc}
\usepackage[utf8]{inputenc}
\usepackage{lmodern}
\usepackage{booktabs}
\usepackage{array}
\usepackage{longtable}
\usepackage{tabularx}
\usepackage{enumitem}
\usepackage{hyperref}
\usepackage{xcolor}

\hypersetup{
  colorlinks=true,
  linkcolor=blue,
  urlcolor=blue
}

\setlength{\parindent}{0pt}
\setlength{\parskip}{0.5em}
\setlist[itemize]{leftmargin=1.4em,itemsep=0.3em,topsep=0.3em}
\setlist[enumerate]{leftmargin=1.6em,itemsep=0.3em,topsep=0.3em}
\renewcommand{\arraystretch}{1.15}

\title{CoT Loop Detector Prefill Follow-up\\Detailed Summary of Reopened Thread Rounds 1--12}
\author{}
\date{2026-03-13 21:39 UTC}

\begin{document}
\maketitle

\section*{Executive Summary}
The reopened thread started from a real concern: the completion-side classifier was only a proof of concept, while the actual goal is \emph{prefill-stage} classification, and a classifier with roughly 50\% precision and 50\% recall is not useful. The work in this thread therefore attacked the problem in three phases: first isolate whether the prefill signal survives stronger confound controls, then test whether representation or training changes can improve the controlled prefill detector, and finally change the \emph{target} itself once the old one-bit label looked exhausted.

The main scientific conclusions are:
\begin{itemize}
  \item The prefill detector is \emph{not} pure noise. After controlling for source and prompt length, and after training against an explicit metadata baseline, hidden-state probes still beat metadata by a nontrivial margin.
  \item Much of the naive prefill lift was shortcut-driven. Prompt length was the dominant confound uncovered in the early rounds.
  \item Representation and objective tuning on the old one-bit target largely saturated. The strongest controlled prefill arm remained the Round 6 metadata-aware all-layer last-token anchor with matched PR-AUC about \texttt{0.6498}.
  \item The major bottleneck was the label itself. On the preserved prompt subsets, small-horizon onset targets were effectively absent: \texttt{loop\_by\_64} was all-negative, and a scout showed zero events through \texttt{1024} and only \texttt{1/256} by \texttt{2048}.
  \item Late-horizon onset is not absent, just late. A long-horizon scout found \texttt{39/256} positives by \texttt{16384} and \texttt{62/256} by \texttt{30000}, with median first hit at token \texttt{11757}.
  \item The next scientifically correct branch is therefore preserved-subset training on \texttt{loop\_by\_16384}, not an immediate composition shift. That round was prepared but not launched because the GPU node was returning devices already occupied by another user's external jobs.
\end{itemize}

\section*{Metric Guide and Jargon}
\begin{description}[style=nextline,leftmargin=1.8em]
  \item[Prefill] Activations computed while the model reads the prompt, before any new answer tokens are generated.
  \item[Completion or posterior] Activations computed after generation begins. These are a useful upper bound, but they are not the target problem here.
  \item[Probe] A small classifier trained on frozen activations. In this project the probes were mainly linear models or multi-layer perceptrons (MLPs).
  \item[PR-AUC] Area under the precision-recall curve. This is usually the better headline metric when positives are rare, because accuracy and even ROC-AUC can look overly flattering under imbalance.
  \item[ROC-AUC] Area under the receiver-operating-characteristic curve. Useful, but not sufficient by itself under imbalance.
  \item[Macro-F1] The F1 score computed for each class and then averaged, so the majority negative class does not dominate the number.
  \item[Metadata baseline] A classifier that only sees cheap prompt-level metadata such as source and prompt length, not hidden states. Here the validated legacy baseline used source, prompt length, newline count, digit count, and LaTeX-dollar count.
  \item[Natural test] The ordinary preserved test distribution. In these rounds it was the \texttt{560}-example MATH-500 plus AIME test set with \texttt{115} positives.
  \item[Matched test] A cleaner evaluation slice where each positive is paired with a source-matched and near-length-matched negative, so prompt-length shortcuts are harder to exploit.
  \item[Source-control] A training subset with the same source and class counts as a stricter control arm, but without matching prompt lengths. This isolates source composition from prompt length.
  \item[Anchor] The main prefill representation used in the controlled rounds: the all-layer last-token summary at the prompt boundary.
  \item[\texttt{loop\_by\_H}] A binary label that is positive if the loop detector fires by generation horizon \texttt{H}. Example: \texttt{loop\_by\_16384} asks whether a loop appears by token \texttt{16384}.
  \item[Onset CDF] The cumulative fraction of prompts whose first loop event has appeared by each horizon. It answers ``how early do loops start?'' rather than only ``do they ever happen?''
  \item[Survival scout] A one-pass rollout used to estimate first-hit timing and right-censoring. Right-censoring means some prompts still have no observed loop by the maximum decode length, so their true onset could be later than the scout limit.
  \item[Within source$\times$length bins] Metrics computed inside groups that share source and coarse prompt-length range, then averaged. This is a harsher check against shortcut exploitation.
\end{description}

\section*{Baseline Before the Reopened Thread}
The reopened loop did \emph{not} start from zero. The relevant pre-thread baseline was Round G from the same project:
\begin{itemize}
  \item Best prefill arm: all-layer prefill last-token concat with macro-F1 \texttt{0.6641 +/- 0.0172}, PR-AUC \texttt{0.4737 +/- 0.0042}, ROC-AUC \texttt{0.7520 +/- 0.0017}, and positive F1 \texttt{0.4955 +/- 0.0129}.
  \item Completion-side upper bound: completion-view probes were much stronger, with PR-AUC around \texttt{0.805} and macro-F1 around \texttt{0.80}.
\end{itemize}

That baseline framed the entire reopened thread. The question was not whether activations contain \emph{any} signal. The question was whether \emph{prefill-only} activations contain enough robust signal once prompt shortcuts are controlled, and if not, whether the right fix is the representation, the training objective, or the label.

\section*{Round-by-Round Spine}
\small
\begin{longtable}{@{}p{0.8cm}p{3.0cm}p{4.3cm}p{6.0cm}@{}}
\toprule
Round & Main question & What changed & Main result and decision \\
\midrule
\endfirsthead
\toprule
Round & Main question & What changed & Main result and decision \\
\midrule
\endhead
1 & Is prefill lift mostly dataset composition? & Dropped GSM8K, then matched source and prompt length on the preserved prefill tensors. & Prefill signal dropped materially under source+length matching but stayed above metadata. Decision: decompose train size, source mix, and prompt length separately. \\
2 & Which confound matters most: train size, source mix, or length? & Added size-control and source-control arms plus a matched evaluation slice. & Prompt length was the dominant shortcut. Decision: stop more broad composition sweeps and train against metadata directly. \\
3 & Does hidden-state signal survive explicit metadata conditioning? & Fit metadata-only baseline first, then trained the hidden probe as an additive residual over the frozen metadata logit. & Yes. Controlled prefill signal remained clearly above metadata. Decision: hold this frame fixed and test better prefill summaries. \\
4 & Can cheap prompt-wide summaries replace the anchor? & Rebuilt fixed splits with multiple new prompt summaries, no new rollouts. & No replacement beat the metadata-aware anchor. Decision: test whether those views help only when fused with the anchor. \\
5 & Is boundary context complementary to the anchor? & Added fusion arms such as \texttt{anchor + last16}. & Modest local lift for \texttt{anchor + last16}. Decision: verify under a cleaner selector and capacity control. \\
6 & Does the Round 5 gain survive a proper holdout selector? & Moved selection to a train-side holdout and added a low-rank fusion control. & No. The anchor became best again. Decision: test whether the bottleneck is objective rather than representation fusion. \\
7 & Can length-robust losses rescue the anchor? & Tried larger-batch BCE, equal-bin BCE, and equal-bin BCE plus within-bin ranking. & No objective tweak beat the Round 6 anchor. Decision: stop tuning the same one-bit anchor classifier and change target or representation. \\
8 & Is one last order-sensitive suffix representation worth trying before changing targets? & Added a strided tail-token suffix view and hybrid anchor+tail arm. & Suffix order helped locally, but no arm beat the earlier Round 6 anchor. Decision: move to richer supervision. \\
9 & Does the smallest onset target work? & Added \texttt{loop\_by\_horizon} labels and rebuilt the preserved fixed subsets at \texttt{H=64}. & \texttt{loop\_by\_64} was all-negative on both preserved subsets. Decision: do not guess another horizon blindly; run a scout. \\
10 & How early do loops actually start on the preserved subset? & Ran a \texttt{256}-prompt scout to \texttt{2048}. & Zero events through \texttt{1024}, only \texttt{1/256} by \texttt{2048}. Decision: small-horizon binaries are dead; run a long-horizon survival scout. \\
11 & Is there enough late-onset mass to stay on the preserved subset? & Ran a \texttt{30000}-token survival scout on the same \texttt{256} prompts. & Yes. \texttt{39/256} by \texttt{16384}, \texttt{62/256} by \texttt{30000}. Decision: stay on preserved split and train \texttt{loop\_by\_16384}. \\
12 & Can the selected late-horizon round be launched immediately? & Prepared the full preserved-subset rebuild and anchor linear-vs-MLP follow-up. & Scientifically ready, operationally blocked. Slurm kept allocating GPUs already occupied by another user's external jobs. \\
\bottomrule
\end{longtable}
\normalsize

\section*{Detailed Chronology}

\subsection*{Round 1: Composition Controls}
\textbf{Question.} If the prefill classifier looks decent, is it because the hidden states truly encode loop risk, or because the probe is exploiting easy prompt-level structure such as source mix and prompt length?

\textbf{What changed.} I reused the preserved Round G prefill tensors, so no new rollouts were needed. First I measured how strong a metadata-only classifier already was on the same natural test set. Then I built two controlled training subsets:
\begin{enumerate}
  \item a \textbf{math-only} subset that drops GSM8K but keeps competition-math and OpenR1-Math;
  \item a \textbf{math-only source+length-matched} subset that pairs positives and negatives by source and prompt-length order.
\end{enumerate}

\textbf{Key numbers.}
\begin{itemize}
  \item Metadata-only baseline on the preserved natural test: ROC-AUC \texttt{0.6590}, PR-AUC \texttt{0.3360}.
  \item Math-only MLP: macro-F1 \texttt{0.6782 +/- 0.0086}, PR-AUC \texttt{0.4704 +/- 0.0302}, ROC-AUC \texttt{0.7482 +/- 0.0060}.
  \item Math-only source+length-matched MLP: macro-F1 \texttt{0.6335 +/- 0.0031}, PR-AUC \texttt{0.4092 +/- 0.0074}, ROC-AUC \texttt{0.7172 +/- 0.0043}.
\end{itemize}

\textbf{Interpretation.} The prefill detector was not just random guessing, but a substantial part of its apparent lift came from coarse prompt structure. Dropping GSM8K alone barely hurt the best prefill arm, which meant the effect was not a single-source artifact. Matching by source and prompt length, however, produced a clear drop. The remaining signal above metadata was real, but clearly smaller than the raw Round G number suggested.

\textbf{Design decision.} Do not spend the next round on larger probes over the same feature family. Instead, separate train-size loss, source reweighting, and prompt-length control more cleanly.

\subsection*{Round 2: Size, Source, and Length Decomposition}
\textbf{Question.} Was the Round 1 loss mainly caused by fewer training examples, by source rebalance, or by prompt-length control?

\textbf{What changed.} I kept the preserved natural test set fixed, reused the same prefill tensors, and added two new train arms:
\begin{itemize}
  \item \textbf{size-control}: same train size as the matched arm, sampled randomly;
  \item \textbf{source-control}: same train size and same per-source/per-class counts as the matched arm, but random over prompt lengths.
\end{itemize}
I also built a new \textbf{matched test} with \texttt{115} positives and \texttt{115} near-length negatives within source, so the evaluation itself was less shortcut-friendly.

\textbf{Key numbers.}
\begin{itemize}
  \item Natural-test hidden PR-AUC:
    \texttt{0.4704} (math-only),
    \texttt{0.4579} (size-control),
    \texttt{0.4562} (source-control),
    \texttt{0.4092} (length-matched).
  \item Matched-test hidden PR-AUC:
    \texttt{0.7020} (math-only),
    \texttt{0.6605} (size-control),
    \texttt{0.6780} (source-control),
    \texttt{0.6211} (length-matched).
  \item Even the strict length-matched arm still beat metadata on the matched test by about \texttt{+0.0829} PR-AUC and \texttt{+0.0968} ROC-AUC.
\end{itemize}

\textbf{Interpretation.} The dominant confound was prompt length. The drop from \texttt{math-only} to \texttt{size-control} was small, so train size was not the main story. The gap from \texttt{size-control} to \texttt{source-control} was also small, so source rebalance by itself was not the main story either. The real collapse happened from \texttt{source-control} to \texttt{length-matched}. This was the point where the thread stopped treating source and length as vague nuisances and started treating prompt length as the central shortcut.

\textbf{Design decision.} Hold the controlled subsets fixed and ask a sharper question: does hidden-state signal survive after \emph{explicitly} conditioning on metadata, not just after changing dataset composition?

\subsection*{Round 3: Residual-over-Metadata Training}
\textbf{Question.} Once a metadata model is already in the prediction, does a prefill hidden-state probe still add useful information?

\textbf{What changed.} I added a dedicated trainer that first reconstructs prompt metadata from preserved sample IDs, fits a metadata-only model, freezes its logit, and then trains the hidden-state probe as an \emph{additive residual} on top of that logit. A smoke run showed that the older ``legacy'' metadata baseline reproduced Round 2 cleanly, while a supposedly stronger interaction-rich baseline actually underperformed, so I kept the validated legacy control instead of overclaiming.

\textbf{Key numbers.}
\begin{itemize}
  \item Legacy metadata baseline:
    natural PR-AUC \texttt{0.3346}, ROC-AUC \texttt{0.6567};
    matched PR-AUC \texttt{0.5487}, ROC-AUC \texttt{0.5408}.
  \item Metadata-residual hidden probe on \texttt{sourcectrl} training:
    natural PR-AUC \texttt{0.4201 +/- 0.0069}, ROC-AUC \texttt{0.7504 +/- 0.0014};
    matched PR-AUC \texttt{0.6474 +/- 0.0088}, ROC-AUC \texttt{0.6857 +/- 0.0073}.
  \item Incremental lift over metadata:
    \texttt{+0.0855} natural PR-AUC and \texttt{+0.0987} matched PR-AUC.
\end{itemize}

\textbf{Interpretation.} This was the first round that made the ``prefill signal is real'' claim substantially harder to dismiss. The hidden-state branch was no longer benefiting merely from a cleaner subset; it had to add log-odds on top of a frozen metadata predictor. It still did. That meant the project had a defensible controlled prefill signal, even if it remained much weaker than completion-side classification.

\textbf{Design decision.} Keep the metadata-aware frame fixed. The next lever should be better prefill summaries, not more train/test composition churn.

\subsection*{Round 4: Prompt-Wide Summary Replacements}
\textbf{Question.} Can a better prompt summary replace the all-layer last-token anchor?

\textbf{What changed.} I rebuilt prefill views on the exact preserved sample IDs and labels, again without new rollouts. New prompt summaries included last-16 pooling, short boundary deltas, late-vs-middle drift, and a small hybrid summary.

\textbf{Key numbers.}
\begin{itemize}
  \item Anchor:
    natural PR-AUC \texttt{0.4182 +/- 0.0114},
    matched PR-AUC \texttt{0.6416 +/- 0.0167},
    matched within-bin PR-AUC \texttt{0.6494}.
  \item Best replacement (\texttt{boundary\_delta8}):
    natural PR-AUC \texttt{0.4099},
    matched PR-AUC \texttt{0.6253}.
  \item Worst replacement (\texttt{late\_vs\_middle\_delta}):
    natural PR-AUC \texttt{0.3702},
    matched PR-AUC \texttt{0.6092}.
\end{itemize}

\textbf{Interpretation.} Cheap prompt-wide summaries were not useless, but none of them beat the anchor as a replacement. The project therefore stopped viewing ``simple pooling over more prompt positions'' as the likely fix.

\textbf{Design decision.} Test whether those views are only helpful as augmentations \emph{on top of} the anchor.

\subsection*{Round 5: Anchor Augmentation}
\textbf{Question.} Is extra boundary context complementary when the anchor remains in the model?

\textbf{What changed.} I extended the residual trainer to concatenate multiple feature views while verifying sample-ID and label alignment, then trained three augmentation arms on the fixed Round 4 tensors.

\textbf{Key numbers.}
\begin{itemize}
  \item Best arm, \texttt{anchor\_plus\_last16}:
    natural PR-AUC \texttt{0.4313 +/- 0.0030},
    matched PR-AUC \texttt{0.6501 +/- 0.0020},
    matched within-bin PR-AUC \texttt{0.6533}.
  \item Relative to the Round 4 anchor:
    \texttt{+0.0130} natural PR-AUC,
    \texttt{+0.0085} matched PR-AUC,
    \texttt{+0.0039} matched within-bin PR-AUC.
  \item Larger raw fusion was not automatically better:
    \texttt{anchor\_plus\_last16\_delta8} lost ground on matched metrics.
\end{itemize}

\textbf{Interpretation.} Round 4's failure was specifically about \emph{replacement}, not about the whole boundary-context family. There was a real but modest gain from adding a coarse last-16 summary on top of the anchor.

\textbf{Design decision.} Before accepting that gain, stress-test it under a proper train-side holdout selector and a capacity-controlled fusion arm.

\subsection*{Round 6: Holdout-Selected Capacity Diagnostic}
\textbf{Question.} Does the Round 5 augmentation gain survive once checkpoint selection is moved off the evaluation set?

\textbf{What changed.} I added a fixed train-side holdout and a low-rank bottleneck option, then ran three arms under the same metadata-aware frame:
\texttt{anchor},
\texttt{anchor\_plus\_last16},
and \texttt{anchor\_plus\_last16\_lowrank256}.

\textbf{Key numbers.}
\begin{itemize}
  \item Anchor:
    natural PR-AUC \texttt{0.4138 +/- 0.0045},
    matched PR-AUC \texttt{0.6498 +/- 0.0052},
    matched within-bin PR-AUC \texttt{0.6641}.
  \item Raw \texttt{anchor + last16}:
    natural PR-AUC \texttt{0.4176 +/- 0.0014},
    matched PR-AUC \texttt{0.6418 +/- 0.0025},
    matched within-bin PR-AUC \texttt{0.6348}.
  \item Low-rank fusion:
    natural PR-AUC \texttt{0.4151 +/- 0.0054},
    matched PR-AUC \texttt{0.6453 +/- 0.0061},
    matched within-bin PR-AUC \texttt{0.6521}.
\end{itemize}

\textbf{Interpretation.} Under cleaner selection, the anchor became best again. The low-rank arm showed that the extra boundary view was not empty, but the apparent Round 5 gain was not robust enough to displace the anchor.

\textbf{Design decision.} Treat fusion as second-order. Next test whether the real bottleneck is the loss function rather than the representation.

\subsection*{Round 7: Objective and Selector Sweep}
\textbf{Question.} Can length-robust objectives or within-bin ranking rescue the anchor without changing features?

\textbf{What changed.} I extended the trainer to support checkpoint selection on within-bin metrics, equal-bin binary cross-entropy (BCE), and an optional within-bin pairwise ranking auxiliary. I also recovered a ``free'' diagnostic by reselecting the saved Round 6 anchor logs using within-bin PR-AUC, so I could isolate selector effects without retraining.

\textbf{Key numbers.}
\begin{itemize}
  \item Free selector-only reselection of the Round 6 anchor:
    matched PR-AUC fell from \texttt{0.6498} to \texttt{0.6254};
    matched within-bin PR-AUC fell from \texttt{0.6641} to \texttt{0.6374}.
  \item Large-batch BCE control:
    matched PR-AUC \texttt{0.6381 +/- 0.0182},
    matched within-bin PR-AUC \texttt{0.6496}.
  \item Equal-bin BCE:
    matched PR-AUC \texttt{0.6189 +/- 0.0058},
    matched within-bin PR-AUC \texttt{0.6146}.
  \item Equal-bin BCE plus ranking:
    matched PR-AUC \texttt{0.6270 +/- 0.0135},
    matched within-bin PR-AUC \texttt{0.6242}.
\end{itemize}

\textbf{Interpretation.} Nothing in this family beat the original Round 6 anchor. Selector alignment was not a free fix, equal-bin BCE overcorrected, and the ranking auxiliary only partly repaired the damage. This was the point where the thread concluded that small training-objective tweaks were not the dominant bottleneck.

\textbf{Design decision.} Stop tuning the same one-bit static-anchor classifier. The next round should change the target or, at most, try one last genuinely order-sensitive representation.

\subsection*{Round 8: Order-Sensitive Tail Representation}
\textbf{Question.} Before paying the full cost of new labels, is one last suffix-order representation check worth doing?

\textbf{What changed.} I confirmed from the builder path that preserved prefill datasets do not store rollout trajectories, so target changes would require new label-build passes rather than cheap relabeling. Given that constraint, I ran one last fixed-label representation test: a \texttt{tail64\_strided\_concat} suffix view that concatenates final-layer prompt-token states at offsets \texttt{\{64, 48, 32, 16, 8, 4, 2, 1\}} plus a short-prompt availability mask. After fixing an extractor dispatch bug and a matched-split loader bug, the three-arm batch completed.

\textbf{Key numbers.}
\begin{itemize}
  \item Rebuilt anchor:
    natural PR-AUC \texttt{0.3957},
    matched PR-AUC \texttt{0.6121},
    matched within-bin PR-AUC \texttt{0.6198}.
  \item \texttt{tail64\_strided\_final}:
    natural PR-AUC \texttt{0.4064},
    matched PR-AUC \texttt{0.6175},
    matched within-bin PR-AUC \texttt{0.6420}.
  \item \texttt{anchor\_plus\_tail64\_lowrank256}:
    natural PR-AUC \texttt{0.4057},
    matched PR-AUC \texttt{0.6253},
    matched within-bin PR-AUC \texttt{0.6307}.
  \item External reference still stronger:
    reproduced Round 6 anchor matched PR-AUC \texttt{0.6498},
    matched within-bin PR-AUC \texttt{0.6641}.
\end{itemize}

\textbf{Interpretation.} Preserving suffix order was a real local gain over the rebuilt Round 8 anchor, especially on within-bin metrics, so the view was not empty. But none of the Round 8 arms beat the earlier controlled Round 6 anchor. That made representation-only work on the old one-bit target look saturated.

\textbf{Design decision.} Move to \emph{richer supervision}. The next round should change the target, starting with compact onset-horizon labels.

\subsection*{Round 9: First Target Change, \texttt{loop\_by\_64}}
\textbf{Question.} Does the smallest meaningful onset label already give a trainable target on the preserved subsets?

\textbf{What changed.} I added a new label surface to the existing stack:
\begin{itemize}
  \item \texttt{eventual\_loop}: the legacy one-bit label;
  \item \texttt{loop\_by\_horizon}: positive if the first loop appears by horizon \texttt{H}.
\end{itemize}
I also added tooling to materialize fixed-split JSONL prompt subsets while preserving original sample IDs, so metadata-aware evaluation stayed aligned to the preserved pools. After a runtime-environment repair and a vLLM memory-envelope repair, the \texttt{loop\_by\_64} batch completed its dataset build.

\textbf{Key numbers.}
\begin{itemize}
  \item Source-control train: \texttt{2096} negatives, \texttt{0} positives.
  \item Source-control test: \texttt{560} negatives, \texttt{0} positives.
  \item Matched train placeholder: \texttt{230} negatives, \texttt{0} positives.
  \item Matched test: \texttt{230} negatives, \texttt{0} positives.
\end{itemize}

\textbf{Interpretation.} This was not a weak training result. It was a stronger statement: on the preserved prompt subsets, \texttt{loop\_by\_64} simply does not occur under the current loop detector. The old plan of trying \texttt{64}, then \texttt{128}, then \texttt{256} as a blind sweep no longer made sense.

\textbf{Design decision.} Stop guessing horizons. Run a scout that measures actual onset prevalence directly.

\subsection*{Round 10: Short-Horizon Scout}
\textbf{Question.} If \texttt{H=64} is empty, where does onset mass actually begin?

\textbf{What changed.} I wrote a scout helper that samples a source-by-length stratified slice, runs one rollout pass, records first-hit timing, and derives post-hoc prevalence for multiple horizons from the same run. I then ran it on a stratified \texttt{256}-prompt source-control-train slice out to \texttt{2048}.

\textbf{Key numbers.}
\begin{itemize}
  \item \texttt{loop\_by\_64}: \texttt{0/256}
  \item \texttt{loop\_by\_128}: \texttt{0/256}
  \item \texttt{loop\_by\_256}: \texttt{0/256}
  \item \texttt{loop\_by\_512}: \texttt{0/256}
  \item \texttt{loop\_by\_1024}: \texttt{0/256}
  \item \texttt{loop\_by\_2048}: \texttt{1/256}
  \item Only observed first-hit step: \texttt{1423}
\end{itemize}

\textbf{Interpretation.} Compact onset binaries were not merely poor at \texttt{64}; they were essentially absent all the way through \texttt{1024}, and still too sparse at \texttt{2048} to justify another full binary rebuild.

\textbf{Design decision.} Shift from binary-onset thinking to a long-horizon survival-style scout. The new question was whether the preserved subset had \emph{late} trainable onset mass or whether composition had to change entirely.

\subsection*{Round 11: Long-Horizon Survival Scout}
\textbf{Question.} Is there enough late-onset mass on the preserved subset to stay on the same fixed prompts, or must the project switch to an onset-enriched composition?

\textbf{What changed.} I ran the same \texttt{256}-prompt stratified preserved slice out to \texttt{30000} tokens with a lighter rollout scheduler, and I consulted Athena repeatedly to turn the result into an explicit branch rule instead of an ad hoc judgment. While the scout was running, I also audited the preserved fixed-subset manifests and confirmed a subtle but important fact: the preserved subset remained balanced on the legacy \texttt{eventual\_loop} bit, yet collapsed to all-negative at \texttt{loop\_by\_64}. So the problem was specifically \emph{late onset}, not generic imbalance.

\textbf{Key numbers.}
\begin{itemize}
  \item \texttt{loop\_by\_2048}: \texttt{2/256} (\texttt{0.78\%})
  \item \texttt{loop\_by\_4096}: \texttt{3/256} (\texttt{1.17\%})
  \item \texttt{loop\_by\_8192}: \texttt{17/256} (\texttt{6.64\%})
  \item \texttt{loop\_by\_16384}: \texttt{39/256} (\texttt{15.23\%})
  \item \texttt{loop\_by\_30000}: \texttt{62/256} (\texttt{24.22\%})
  \item First-hit timing among positives:
    minimum \texttt{810},
    p25 \texttt{7209},
    median \texttt{11757},
    p75 \texttt{19552},
    p90 \texttt{23223},
    maximum \texttt{28403}.
\end{itemize}

\textbf{Interpretation.} This round resolved the main scientific fork. Onset was still essentially absent at small and medium horizons, but there \emph{was} enough late-onset mass by \texttt{16384} to support a preserved-subset follow-up. The earliest viable late-horizon target was therefore \texttt{loop\_by\_16384}. Going all the way to \texttt{30000} would have pushed the label too deep into the right tail; staying at \texttt{8192} would still have been too sparse.

\textbf{Expected full-subset counts if the scout prevalence carries through.}
\begin{itemize}
  \item Source-control train: about \texttt{319} positives.
  \item Source-control test: about \texttt{85} positives.
  \item Matched test: about \texttt{35} positives.
\end{itemize}

\textbf{Design decision.} Stay on the preserved split and train a coarse late-horizon target at \texttt{loop\_by\_16384}. Do \emph{not} pivot immediately to an onset-enriched composition shift.

\subsection*{Round 12: Prepared but Resource-Blocked}
\textbf{Question.} Can the selected \texttt{loop\_by\_16384} branch be executed immediately?

\textbf{What changed.} I prepared a follow-up batch that rebuilds the preserved source-control and matched datasets using \texttt{label\_target=loop\_by\_horizon@16384}, then trains the anchor representation with a small linear-versus-MLP comparison under the same preserved labels.

\textbf{What was ready.}
\begin{itemize}
  \item New launcher:
    \texttt{round12\_h16384\_target.sbatch}
  \item Safer rollout scheduler:
    \texttt{dp=4},
    \texttt{max\_num\_seqs=8},
    \texttt{max\_num\_batched\_tokens=4096},
    \texttt{VLLM\_GPU\_MEMORY\_UTILIZATION=0.20}
  \item Scientific target fixed:
    preserve the same prompt subsets and train \texttt{loop\_by\_16384}
\end{itemize}

\textbf{Why it did not launch.} The GPU node never became safe enough. Slurm looked empty from its own queue view, but \texttt{nvidia-smi} still showed another user's large external Python jobs sitting on GPUs \texttt{2}, \texttt{4}, and \texttt{6}. A fresh \texttt{--gres=gpu:4} probe kept returning \texttt{2,3,4,6}, which overlaps those occupied devices. Launching the real round into that state would very likely have produced a CUDA out-of-memory failure rather than a scientific result.

\textbf{Interpretation.} The blocker at the end of the thread was operational, not scientific. The design question was settled; the scheduler environment was not.

\section*{What the Thread Established}
\begin{enumerate}
  \item \textbf{Controlled prefill signal exists.} After explicit source and prompt-length controls, and after direct conditioning on metadata, prefill hidden states still contain loop-relevant information.
  \item \textbf{Shortcut pressure is real and large.} Prompt length was the dominant confound in the early rounds, and many seemingly promising gains weakened once the evaluation became more controlled.
  \item \textbf{Round 6 is the best controlled old-target baseline.} The metadata-aware anchor with train-side holdout selection remains the strongest defensible prefill detector on the legacy label, with matched PR-AUC around \texttt{0.6498}.
  \item \textbf{Minor representation and objective tweaks are largely exhausted on the old one-bit target.} Boundary-context fusion, low-rank fusion, equal-bin losses, and within-bin ranking all failed to produce a new best model.
  \item \textbf{The real label problem is timing.} Small-horizon onset labels are almost absent on the preserved prompt sets, but late-horizon onset is not.
  \item \textbf{The earliest viable next target is \texttt{loop\_by\_16384}.} That is the central output of the entire Round 9--11 sequence.
\end{enumerate}

\section*{Recommended Immediate Next Step}
If the node clears, submit the already prepared preserved-subset Round 12 batch:
\begin{itemize}
  \item rebuild preserved \texttt{sourcectrl} and \texttt{matched} datasets at \texttt{loop\_by\_16384};
  \item train the anchor view with a small linear-versus-MLP comparison;
  \item evaluate matched performance primarily with PR-AUC and within-bin metrics, not accuracy alone.
\end{itemize}

If the node remains contested, the next intervention should be infrastructural rather than scientific: reserve or isolate safe GPUs, move to a cleaner node, or otherwise avoid Slurm allocations that overlap with external non-Slurm jobs. The thread already answered the design question it was supposed to answer.

\section*{Bottom Line}
The reopened loop did not prove that prefill-stage loop detection is solved. It did something more important: it separated a misleading short-horizon story from a viable late-horizon story. The prefill detector is weaker than the completion detector, and much of the naive prefill lift was shortcut-inflated. But after controlling for those shortcuts, the signal is still real, and the long-horizon scout shows that the preserved prompt subsets are not hopeless. They are simply \emph{late}. The scientifically correct next experiment is therefore the preserved \texttt{loop\_by\_16384} classifier round, and the reason it is not already finished is GPU allocation safety, not lack of a coherent target.

\end{document}
